\documentclass[11pt]{article}
\usepackage[margin=1in]{geometry}
\usepackage{amsmath,amssymb}
\usepackage{booktabs}
\usepackage{graphicx}
\usepackage{hyperref}
\usepackage{natbib}
\usepackage{algorithm}
\usepackage{algpseudocode}

\title{NEXUS: More Agents Does Not Mean Better Performance \\
in Multi-Agent Software Repair}

\author{
Kaushal T. \\
\texttt{github.com/Kaushalt2004/cognicore-my-openenv}
}

\date{May 2026}

\begin{document}
\maketitle

\begin{abstract}
We present NEXUS, a multi-agent orchestration runtime for autonomous software repair
that coordinates specialized AI agents (Planner, Coder, Reviewer, Tester, Verifier)
with persistent cross-session memory and learned routing policies.
Through systematic ablation across 4 routing policies, 20 SWE-bench-lite tasks,
and 280+ recorded trajectories, we demonstrate a counterintuitive finding:
\textbf{adding a code review agent reduces solve rate by 5\% while increasing
token cost by 35\%}. We show that test-before-review ordering outperforms
review-before-test, that persistent episodic memory accumulates without
performance degradation over long sessions, and that the optimal
``cognition-per-token'' ratio favors minimal agent pipelines.
Our trajectory store enables offline reinforcement learning over
coordination policies, framing multi-agent software engineering
as a sequential decision-making problem.
NEXUS is open-source and integrates with LangChain, Gemini, GPT-4, and Claude.
\end{abstract}

\section{Introduction}

The dominant paradigm in autonomous software engineering assumes that
more sophisticated multi-agent architectures yield better results.
Systems like SWE-Agent \citep{yang2024sweagent}, AutoCodeRover \citep{zhang2024autocoderover},
and Devin coordinate multiple specialized agents for planning, coding,
reviewing, and testing. However, the marginal contribution of each
agent type to overall solve rate remains poorly understood.

We introduce NEXUS, a modular multi-agent runtime built on the
CogniCore cognition engine, designed specifically to answer:
\textit{Which agents actually help, and when do they hurt?}

Our key contributions:
\begin{enumerate}
    \item \textbf{Empirical evidence that review agents hurt performance}
    in rule-based pipelines, reducing solve rate from 95\% to 90\%
    while adding 35\% token overhead (Section~\ref{sec:reviewer}).
    \item \textbf{A persistent cognition engine} that accumulates
    cross-session repair memory without performance degradation
    over 3+ consecutive sessions (Section~\ref{sec:memory}).
    \item \textbf{A trajectory store} for offline RL over coordination
    policies, with 280+ recorded trajectories and computed rewards
    (Section~\ref{sec:trajectory}).
    \item \textbf{Token economics analysis} showing the Pareto-optimal
    policy uses only 1,446 tokens per solved task (Section~\ref{sec:economics}).
\end{enumerate}

\section{System Architecture}

NEXUS consists of three layers:

\paragraph{AXIOM Orchestration Layer.}
Routes tasks through agent pipelines via a Coordination Bus.
Four routing policies are supported:
\texttt{minimal} (Coder$\to$Tester),
\texttt{standard} (Memory$\to$Planner$\to$Coder$\to$Reviewer$\to$Tester),
\texttt{test\_first} (Memory$\to$Planner$\to$Coder$\to$Tester$\to$Reviewer),
and \texttt{review\_first} (Memory$\to$Planner$\to$Coder$\to$Reviewer$\to$Coder$\to$Tester).

\paragraph{CogniCore Cognition Engine.}
Provides persistent episodic memory (SQLite-backed), a reflection engine
that generates insights from failure patterns, and semantic patch rejection
using combined edit-distance and structural similarity.

\paragraph{Infrastructure Layer.}
Subprocess-isolated sandbox execution, deterministic event logging for
replay, a token tracker with per-agent cost accounting, and a trajectory
store for offline RL training data.

\section{Experimental Setup}

\subsection{Benchmark}
We evaluate on 20 curated SWE-bench-lite tasks spanning 9 repositories
(Django, SymPy, Flask, Requests, Astropy, scikit-learn, matplotlib, pytest)
across 12 bug categories (arithmetic, encoding, none\_handling, off\_by\_one,
parsing, recursion, type\_conversion, etc.).

\subsection{Agents}
Six agent types: Memory (retrieves past experiences), Planner (selects
repair strategy), Coder (generates patches), Reviewer (validates patch
quality), Tester (executes in sandbox), Verifier (safety checks).

\subsection{Evaluation Protocol}
Each policy runs with 5 maximum attempts per task. We report solve rate,
total attempts, token usage, cost per solve, and per-category breakdown.
Stress tests include multi-seed stability (3 seeds), cold vs. warm memory,
and long-session drift (3 consecutive runs).

\section{Results}

\subsection{Policy Comparison}

\begin{table}[h]
\centering
\caption{NEXUS benchmark results across 4 routing policies.}
\label{tab:policies}
\begin{tabular}{lccccr}
\toprule
\textbf{Policy} & \textbf{Solved} & \textbf{Rate} & \textbf{Tokens} & \textbf{Tok/Solve} & \textbf{Calls} \\
\midrule
test\_first & \textbf{19/20} & \textbf{95.0\%} & 32,854 & 1,729 & 136 \\
minimal & 19/20 & 95.0\% & \textbf{27,476} & \textbf{1,446} & \textbf{68} \\
standard & 18/20 & 90.0\% & 37,118 & 2,062 & 172 \\
review\_first & 18/20 & 90.0\% & 45,591 & 2,533 & 192 \\
\bottomrule
\end{tabular}
\end{table}

\subsection{Reviewer Ablation}
\label{sec:reviewer}

The reviewer agent reduces solve rate from 95\% to 90\% ($-1$ task)
while adding 9,642 tokens (+35\%) and 104 additional agent calls.
Analysis reveals the reviewer's similarity-based rejection mechanism
blocks valid fixes that share structural patterns with the original
buggy code---a fundamental limitation when patches are minor
modifications rather than complete rewrites.

The task exclusively failed by reviewer-containing policies
(\texttt{SWE-dj-15498}) involves a guard fix that is structurally
similar to the buggy code but semantically correct.

\subsection{Persistent Memory}
\label{sec:memory}

\begin{table}[h]
\centering
\caption{Long-session stability: 3 consecutive runs without memory reset.}
\label{tab:drift}
\begin{tabular}{lccc}
\toprule
\textbf{Run} & \textbf{Solved} & \textbf{Attempts} & \textbf{Tokens} \\
\midrule
1 (cold) & 18/20 & 38 & 37,118 \\
2 (warm) & 18/20 & 38 & 37,608 \\
3 (warm) & 18/20 & 38 & 37,608 \\
\bottomrule
\end{tabular}
\end{table}

Memory tokens increase by 490 per session as experiences accumulate,
but this causes no performance degradation or strategy drift.
Zero variance across 3 seeds confirms deterministic behavior.

\subsection{Token Economics}
\label{sec:economics}

The \texttt{minimal} policy achieves the optimal cognition-per-token ratio
at 1,446 tokens per solved task. Adding memory and planning
(\texttt{test\_first}) costs 283 additional tokens per solve (+19.6\%)
with no solve rate improvement, but provides the infrastructure for
LLM-backed agents to leverage memory context in future iterations.

\subsection{LLM Integration}
\label{sec:llm}

With Gemini 2.0 Flash as the coder backend (with rule-based fallback
for rate-limited requests), NEXUS achieves 19/20 (95.0\%) solve rate
with the \texttt{test\_first} policy, confirming that the cognitive
infrastructure works with real LLM backends.

\section{Trajectory Store for Offline RL}
\label{sec:trajectory}

Each benchmark run records full coordination trajectories:
\[
\tau = \{(s_t, a_t, r_t)\}_{t=1}^{T}
\]
where $s_t$ is the task state, $a_t$ is the agent action, and $r_t$ is
the step reward. Trajectory-level reward is:
\[
R(\tau) = r_{\text{solve}} - c_{\text{attempt}} \cdot n - c_{\text{token}} \cdot T/1000
\]

With 280+ trajectories across 4 policies, this enables offline RL training
of coordination policies that learn which agents to invoke given task
characteristics and budget constraints.

\section{Related Work}

SWE-Agent \citep{yang2024sweagent} uses a single agent with tool access.
AutoCodeRover \citep{zhang2024autocoderover} combines search with LLM reasoning.
Neither systematically ablates the contribution of individual agent types
or records trajectories for coordination learning. Our finding that
review agents hurt performance echoes observations in human software
engineering that premature code review can slow development velocity.

\section{Limitations}

\begin{itemize}
    \item Rule-based agents cannot exploit memory hints; the memory pipeline's
    full value requires LLM-backed agents that read and adapt to context.
    \item Our benchmark uses 20 curated tasks; scaling to the full
    SWE-bench (300+ tasks) is ongoing work.
    \item The reviewer finding may not generalize to LLM-backed reviewers
    that perform semantic (not structural) similarity checks.
\end{itemize}

\section{Conclusion}

NEXUS demonstrates that in multi-agent software repair, the relationship
between agent count and performance is non-monotonic. Our key finding---that
review agents reduce solve rate while increasing cost---challenges the
assumption that more agents yield better outcomes. The optimal policy
(\texttt{test\_first}) achieves 95\% solve rate with 1,729 tokens per
solve, providing memory and planning context without reviewer overhead.
Our open-source trajectory store enables the community to train offline
RL coordination policies on real multi-agent repair data.

\bibliographystyle{plainnat}
\begin{thebibliography}{9}
\bibitem[Yang et al.(2024)]{yang2024sweagent}
Yang, J., Jimenez, C.E., Wettig, A., Lieret, K., Yao, S., Narasimhan, K., and Press, O.
\newblock SWE-agent: Agent-Computer Interfaces Enable Automated Software Engineering.
\newblock \emph{arXiv preprint arXiv:2405.15793}, 2024.

\bibitem[Zhang et al.(2024)]{zhang2024autocoderover}
Zhang, Y., Ruan, W., Fan, Z., and Roychoudhury, A.
\newblock AutoCodeRover: Autonomous Program Improvement.
\newblock \emph{arXiv preprint arXiv:2404.05427}, 2024.
\end{thebibliography}

\end{document}
